\documentclass[14pt,a4paper]{extarticle}
\usepackage{fontspec}
\usepackage{polyglossia}
\setmainlanguage{russian}
\setmainfont{Times New Roman}
\usepackage{setspace}
\onehalfspacing
\usepackage{geometry}
\geometry{left=30mm, right=15mm, top=20mm, bottom=20mm}
\usepackage{indentfirst}
\setlength{\parindent}{1.25cm}
\setlength{\parskip}{0pt}
\usepackage{xurl}
\urlstyle{same}
\usepackage{hyperref}
\usepackage{tabularx}
\usepackage{array}
\usepackage[table]{xcolor}
\definecolor{pastelgreen}{HTML}{B5EAD7}
\usepackage{booktabs}
\usepackage{ragged2e}
\usepackage{wasysym}
\usepackage{graphicx}
\usepackage{float}
\usepackage{pdflscape}


\begin{document}

	\sloppy
	\thispagestyle{empty}
	\centering{
		МИНИСТЕРСТВО НАУКИ И ВЫСШЕГО ОБРАЗОВАНИЯ\\
		РОССИЙСКОЙ ФЕДЕРАЦИИ\\
		Федеральное государственное автономное образовательное учреждение высшего образования <<Санкт-Петербургский политехнический университет Петра Великого>>\\[10pt]
		Институт компьютерных наук и кибербезопасности\\
		Высшая школа технологий искусственного интеллекта\\
		02.03.01 Математика и компьютерные науки\\
		\vspace{\fill}
		{\large
			Отчёт по курсовой работе\\
			по дисциплине\\
			{\Large \bfseries Теория графов}\\
			\vspace{\fill}
		}


		{
			\flushright
			{\bfseries Выполнил:}\\
			студент группы 5130201/40003\\
			Джабраилов А.В.\\
			\vspace{2pt}
			\hspace{9cm}Подпись: \hrulefill \\
			\vspace{2pt}
			{\bfseries Принял:}\\
			ст. преп.\\
			Востров А.В.\\
			\vspace{2pt}
			\hspace{9cm}Подпись: \hrulefill \\
			\vspace{3pt}
			\hspace{9cm}<<\rule{1.5cm}{0.4pt}>> \hrulefill\ 2026~г.
		}
		
		\centering{
			Санкт-Петербург, 2026
		}
	}

\newpage
\justifying
\tableofcontents
	
\newpage

\section*{Введение}
\addcontentsline{toc}{section}{Введение}

\section{Математическое описание}

\newpage

\section{Особенности реализации}
	
\newpage
				
\section{Результаты работы программы}
				
\newpage

\section*{Заключение}
\addcontentsline{toc}{section}{Заключение}
				
\paragraph{Плюсы программы.}
				
\paragraph{Минусы программы.}
				
\paragraph{Масштабирование.} 
				
\newpage
				
\begin{thebibliography}{99}
\addcontentsline{toc}{section}{Список литературы}

    %\bibitem{oop} Павловская Т. А., Щюпак Ю. А. С++ Объектно-ориентированное программирование: Практикум. — СПб.: Питер, 2006 — 265 с.

\end{thebibliography}

\newpage
		
% \section*{Приложение 5. Файл utils.h}
% \addcontentsline{toc}{section}{Приложение 5. Файл utils.h}
% \lstinputlisting[language=C++,caption={Файл utils.h.}]{utils.h}
% \newpage

\end{document}